% chapters/journal/01-introduction.tex
% 第1章 引言

\section{引言}

区块链系统中，区块构建者负责从交易池中选取待处理交易并确定其在区块内的执行顺序。以以太坊为代表的智能合约平台上，交易排序质量对区块手续费收益、交易执行成功率和用户体验具有直接影响\cite{daian2019flash}。当前主流排序策略采用基于 Gas Price 的贪心方法，按交易手续费从高到低排列以最大化短期收益。该策略实现简单、计算代价低，但存在明显局限：一方面，贪心排列忽略交易之间的状态读写依赖与资源冲突，当两笔交易访问相同合约状态时，顺序不当可能导致执行失败或回滚；另一方面，低手续费交易在该策略下长期得不到确认，等待时间分布容易失衡。

最大可提取价值（Maximal Extractable Value，MEV）问题进一步加剧了上述矛盾\cite{daian2019flash}。MEV 是指区块构建者通过插入、删除或重排交易所能获得的额外利润，其规模在以太坊生态中已达到相当水平。典型的 MEV 攻击形式包括抢跑交易（front-running）、尾随交易（back-running）和三明治攻击（sandwich attack）\cite{chi2024remeasure}。在三明治攻击中，攻击者在目标交易前后分别插入买入和卖出交易，利用目标交易造成的价格变动获利，而目标交易的用户则承受更大的价格滑点。这类行为使普通用户遭受实际经济损失，同时推高了链上 Gas 费用的整体水平。针对 MEV，业界提出了 MEV-Boost\cite{flashbots2022mevboost}、提议者--构建者分离（PBS）\cite{ethereum2024pbs}、公平排序服务和加密内存池等多种机制。然而，这些方案主要在宏观架构层面重塑 MEV 的分配方式或防止交易信息泄露，在构建者内部如何排序交易以平衡收益与风险，仍缺乏系统的算法设计\cite{alipanahloo2024mev}。

近年来，强化学习（Reinforcement Learning，RL）被引入区块链交易排序领域。Wen 等\cite{wen2025rlbes}提出 RL-BES 系统，将深度强化学习与搜索算法结合用于区块构建；也有研究采用深度 Q 网络（DQN）学习排序策略以抑制 MEV。然而，现有 RL 方法在动作空间设计上存在共性不足：其动作空间通常定义为在 FIFO、Gas 优先等少数预定义规则之间切换，本质上是“选规则”而非“逐笔选交易”。这种粗粒度动作设计将排序决策简化为分类任务，难以捕获交易间细粒度依赖关系，也难以根据候选交易集合构成动态调整策略。

基于上述分析，本文提出一种基于深度强化学习的以太坊交易序列决策排序方法。与已有工作的关键区别在于：本文将区块构建过程建模为逐笔交易选择的马尔可夫决策过程（MDP），智能体在每一步从候选交易集合中选取一笔交易加入区块执行序列，直至满足终止条件。这种建模方式将排序决策粒度从“选择排序规则”细化为“选择具体交易”，使策略能够利用当前位置、剩余容量和已选序列上下文做条件决策。本文主要贡献包括三个方面：（1）构建中高负载区块构建场景下的可配置交易排序仿真框架，统一候选池规模、风险比例、Gas 上限与统计评估设置；（2）提出收益--等待时间均衡--风险暴露三目标统一优化的 RL 排序方法，在奖励中引入增量公平项与位置敏感风险惩罚；（3）在统一实验矩阵下，与 FIFO、Gas 优先、启发式风险感知、Fee-Risk 线性评分和 Fair-Fee 双目标贪心五类基线进行系统对比，并通过鲁棒性与消融分析验证方法稳定性与可解释性。
